% Generated draft: requires mathematical and source review.
\documentclass[11pt]{article}
\usepackage[margin=25mm]{geometry}
\usepackage{fontspec}
\setmainfont{Latin Modern Roman}
\usepackage{amsmath,amssymb,mathtools,unicode-math}
\setmathfont{Latin Modern Math}
\usepackage{enumitem,needspace,xurl,hyperref}
\setlength{\parindent}{0pt}
\setlength{\parskip}{5pt}
\setlength{\emergencystretch}{3em}
\newcommand{\N}{\mathbb{N}}
\newcommand{\Z}{\mathbb{Z}}
\newcommand{\Q}{\mathbb{Q}}
\newcommand{\R}{\mathbb{R}}
\newcommand{\C}{\mathbb{C}}
\newcommand{\F}{\mathbb{F}}
\newcommand{\A}{\mathbb{A}}
\newcommand{\PP}{\mathbb{P}}
\newcommand{\E}{\mathbb{E}}
\newcommand{\eps}{\varepsilon}
\newcommand{\dd}{\,\mathrm{d}}
\begin{document}
\section*{UnsolvedMath ID: 30001333}
\textbf{Source problem code:} catalogue entry
\subsection{Definitions and mathematical statement}
A topological groupoid \(G\rightrightarrows G^{(0)}\) consists of a space \(G\) of arrows, a space \(G^{(0)}\) of objects, continuous source and range maps, and continuous unit, inversion, and composition maps satisfying the groupoid axioms; composition is defined on composable pairs. A topological space is Hausdorff if any two distinct points have disjoint open neighbourhoods. It is locally Hausdorff if each point has an open Hausdorff neighbourhood; this does not require the whole space to be Hausdorff.

For an action of \(G\) on a \(C^*\)-algebra \(A\), write \(A\rtimes G\) schematically for an associated crossed product, and write \(KK^G(A,B)\) schematically for equivariant bivariant \(K\)-theory of suitable coefficient algebras \(A\) and \(B\). These symbols express the objects whose extension is requested. The supplied record does not specify the admissible actions or coefficients, the full or reduced crossed-product convention, or a definition of equivariant theory in the proposed non-Hausdorff setting.

\textbf{Source question.} Explain the role of the Hausdorff hypothesis on \(G\) in the Baum--Connes result for amenable groupoids attributed in the source to Tu. Determine how crossed products and equivariant bivariant \(KK\)-theory can be extended to the non-Hausdorff, locally Hausdorff actions arising in that context, as part of the associated Baum--Connes program.

\textbf{Unresolved scope.} This is a research-program question rather than a fully specified quantified conjecture. The supplied record does not state Tu's precise theorem, its remaining hypotheses, the intended notion of amenability, or the exact meaning of the resulting actions. In particular, it does not identify whether local Hausdorffness is a condition on the arrow space, an action space, or another object in the action formalism. Thus no universal extension theorem, precise assembly-map assertion, or fixed class of coefficients can be extracted without further source clarification. The topological definitions above do not impose additional hypotheses on the intended theorem.
\subsection{Short English statement}
The question asks why Hausdorffness is needed in Tu's cited Baum--Connes result for amenable groupoids, and how crossed products and equivariant \(KK\)-theory should extend to the non-Hausdorff, locally Hausdorff actions under discussion. The supplied record leaves the precise theorem, action formalism, and extension requirements unspecified. Its dated curation labels the program partially solved; this draft does not independently verify that assessment.
\subsection{Sources}
\begin{itemize}
\item Local UnsolvedMath registry, permanent ID \(30001333\), source problem code \texttt{catalogue entry}, title \emph{Baum--Connes Theory for Non-Hausdorff Groupoids}, in the Oberwolfach Reports Open Problems collection. The registry gives proposed year 2009 and does not identify a proposer. Its statement attributes the relevant amenable-groupoid result to Tu, without supplying a separate bibliographic reference for that theorem.
\item Supplied authoritative source citation: \emph{Noncommutative Geometry (2010)}, Oberwolfach Reports, pp.\ 2271--2334, DOI \texttt{10.4171/owr/2009/41}, \url{https://doi.org/10.4171/owr/2009/41}. The literature-source label is ``original Oberwolfach report 4084.'' The supplied background describes these reports as documenting mathematical workshops and their open-problem discussions. The relationship between the citation's 2010 label and the record's 2009 proposed year is not explained.
\item Supplied literature source, labelled ``Buss-Meyer-Zhu non-Hausdorff symmetries'': \url{https://arxiv.org/abs/0907.0339}.
\item Supplied literature source, labelled ``higher-category actions'': \url{https://arxiv.org/abs/0908.0455}.
\item Literature assessment dated exactly \texttt{2026-08-21}: the supplied curation labels the problem \texttt{partially\_solved} and reports that later primary work defines crossed-module and higher-category actions for non-Hausdorff symmetries, addressing the obstruction in the question, while the full crossed-product and Baum--Connes program stated in the report is not completely closed. This is supplied dated curation, not independently verified here; the supplied background explicitly calls for independent verification of time-sensitive status and cited claims.
\item The registry labels the statement \texttt{corrected\_verified} and reports that its mathematical content was checked against the official report, with notation, context, and question form normalized without changing the source-level claim. This is a supplied verification assertion, not a source inspection performed for this draft.
\end{itemize}
\end{document}
